\scanpage{91}
\mspage{13}
Not only the condition \unsure{of|of} \deleted{continuity} ``extendability'' --- it is better to say of inducibility'' --- in the form specified, or weakened (with $\Theta$ replaced by a $\Theta^n$), one finds an additional condition, when one has \inserted{in addition} a morphism
\[
\widehat\Pi_1(C,I)\longrightarrow\widehat\Pi_1(X,I),
\]
that the latter pass to the quotient --- which is expressed very exactly by the condition that for \uline{one} $x_0\in I$, the action of $\Gamma$ on
\[
\widehat\Pi_1(C,x_0)\subset\widehat\Pi_1(U_{0,3},f(x_0))
\]
leave invariant the kernel of
\[
\widehat\Pi_1(C,x_0)\longrightarrow\widehat\Pi_1(X,x_0);
\]
then this will also be true for all the other $x_i\in I$.

It is necessary to express the possibility of lifting the action \deleted{of} $\Theta$ to $\Pi_1(C,I)$. Let us note in this connection

\uline{Lemma} Let $\mathcal C'\unsure{\mathbin{-}|\to}\mathcal C$ be a homomorphism of connected groupoids, which is injective on the $\Pi_1$, \unsure{\deleted{and}|(} so that $\mathcal C'$ is equivalent to a covering of $\mathcal C$, but one does not assume that $\mathcal C'$ is itself a covering of $\mathcal C$, but still that
\margin{90}
